\documentclass[11pt]{amsart}
\usepackage[T1]{fontenc}
\usepackage{lmodern,amsmath,amssymb,amsthm,microtype}
\usepackage[margin=1.15in]{geometry}
\usepackage[colorlinks=true,linkcolor=blue,citecolor=blue,urlcolor=blue]{hyperref}
% Repository locator for the public (non-anonymous) build of
% non_mf_groups_exist.tex.  Kept in a separate file so that the anonymous
% build has an anonymous source archive and not merely an anonymous PDF:
% exclude this file from a double-blind submission and the manuscript still
% compiles, printing "supplied with this submission as supplementary
% material" in place of the link.
%
% Loaded only when \anonymousfalse is set in the preamble.
\newcommand{\coderepository}{https://github.com/SauersML/group-approximation}
\newcommand{\coderepositoryname}{github.com/SauersML/group-approximation}

\hypersetup{pdftitle={An infinite simple Kazhdan sofic group},pdfauthor={\paperauthorname}}
\newtheorem{theorem}{Theorem}
\newtheorem{corollary}[theorem]{Corollary}
\newcommand{\F}{\mathbb F}
\newcommand{\Z}{\mathbb Z}
\newcommand{\LC}{\operatorname{LC}}
\newcommand{\EL}{\operatorname{EL}}
\newcommand{\GL}{\operatorname{GL}}
\newcommand{\PSL}{\operatorname{PSL}}
\newcommand{\doi}[1]{\href{https://doi.org/#1}{\nolinkurl{doi:#1}}}
\raggedbottom
\emergencystretch=1.5em

\begin{document}
\title{An infinite simple Kazhdan sofic group}
\author{\paperauthorname}
\date{}
\subjclass[2020]{Primary 20E32; Secondary 20E26, 20F10, 22D10, 37B10, 16S35}

\begin{abstract}
For every infinite minimal subshift $X$, the group
$\EL_3(\LC(X,\F_2)\rtimes\Z)$ is infinite, simple, finitely generated,
and has property~\textup{(T)}. It is locally embeddable into finite
groups, so it is sofic and hyperlinear. This answers the question of
Brown and Ozawa whether an infinite simple Kazhdan group can be
hyperlinear, and Pestov's sofic version of it. Every Turing degree
occurs as the word-problem degree of one of these groups.
\end{abstract}
\maketitle

Can an infinite simple Kazhdan group be hyperlinear? Brown asked this
in its von Neumann algebra form in 2001~\cite[\S11, Question~7]{Brown},
and Ozawa in the hyperlinear form in 2003~\cite[p.~527]{Ozawa}.
Pestov's Open question~9.1 adds the sofic version~\cite{Pestov}.

\begin{theorem}\label{thm:main}
Let $X\subseteq A^{\Z}$ be an infinite minimal subshift over a finite
alphabet, with shift $T$, and let $\LC(X,\F_2)$ be the ring of locally
constant functions $X\to\F_2$. Then
\[
  G_X=\EL_3\bigl(\LC(X,\F_2)\rtimes_T\Z\bigr)
\]
is an infinite, finitely generated, simple group with Kazhdan's
property~\textup{(T)}. It is locally embeddable into finite groups
\textup{(LEF)}, so it is sofic and hyperlinear.
\end{theorem}

Property~\textup{(T)} follows from the theorem of Ershov and
Jaikin-Zapirain~\cite{EJZ}. The finite models use periodic sequences
with the same short words as $X$, as in Grigorchuk and
Medynets~\cite{GM}. The new step is simplicity. A nontrivial normal
subgroup contains a nontrivial commutator lying in a copy of the finite
simple group $\GL_d(\F_2)$ over a clopen tower, so it contains this
group and with it an elementary matrix. For derived topological full
groups, Matui showed in the same way that a nontrivial normal subgroup
meets a simple union of alternating groups on
towers~\cite[Lemma~3.4 and Theorem~4.9]{Matui}. Thom constructed a
finitely generated Kazhdan LEF group that is not residually
finite~\cite{Thom}, but his example is not simple. A finitely presented
LEF group is residually finite~\cite{VershikGordon}, and $G_X$ is not,
so $G_X$ is not finitely presented. We do not know whether some
finitely presented infinite simple Kazhdan group is sofic.

\section{Proof of Theorem~\ref{thm:main}}

\subsection*{The ring and property (T)}
Write $e_U$ for the indicator of a clopen set $U\subseteq X$. Our
conventions are $(Tx)_t=x_{t+1}$ and $ufu^{-1}=f\circ T^{-1}$ in
$R=\LC(X,\F_2)\rtimes_T\Z=\bigoplus_{j\in\Z}\LC(X,\F_2)\,u^j$. So every
element of $R$ is uniquely a finite sum $\sum_jf_ju^j$, and
$ue_Uu^{-1}=e_{TU}$.

The ring $R$ is generated by $u^{\pm1}$ and the letter indicators
$e_a=e_{\{x:x_0=a\}}$, since products of the translates $u^je_au^{-j}$
are the cylinder indicators, which span $\LC(X,\F_2)$. Let $G=\EL_3(R)$
be the subgroup of $\GL_3(R)$ generated by the elementary matrices
$e_{ij}(r)=I_3+rE_{ij}$, $i\ne j$, $r\in R$, where $E_{ij}$ is the usual
matrix unit. With $[g,h]=ghg^{-1}h^{-1}$,
\begin{equation}\label{eq:elementary}
 \begin{aligned}
 e_{ij}(r+s)&=e_{ij}(r)e_{ij}(s),\qquad i\ne j,\\
 [e_{ik}(r),e_{kj}(s)]&=e_{ij}(rs),\qquad i,j,k\text{ distinct}.
 \end{aligned}
\end{equation}
So the matrices $e_{ij}(s)$ with $s\in\{u,u^{-1}\}\cup\{e_a:a\in A\}$
generate $G$. By Ershov and Jaikin-Zapirain~\cite[Theorem~1.1]{EJZ},
$\EL_n(S)$ has property~\textup{(T)} for every finitely generated
associative unital ring $S$ and every $n\ge3$. So $G$ has
property~\textup{(T)}. It is infinite because $e_{12}(\LC(X,\F_2))$ is
infinite.

\subsection*{Finite models}
A group is LEF~\cite{VershikGordon} if every finite subset embeds
injectively into a finite group, preserving all products that stay in
that subset.

Fix $x\in X$ and $\ell\ge0$. By minimality every word of $X$ occurs in
$x_{[0,\infty)}$, and $x_{[0,2\ell)}$ recurs at arbitrarily large
positions. Choose such a position $m_\ell$ with all words of length
$2\ell+1$ of $X$ occurring in $x_{[0,m_\ell+2\ell)}$. This word begins
and ends with $x_{[0,2\ell)}$, so it has period $m_\ell$, and the
$m_\ell$-periodic sequence $y_\ell$ that agrees with $x$ on $[0,m_\ell)$
has the same words of length $2\ell+1$ as $X$. So $m_\ell$ is at least
the number of these words, which tends to infinity because $X$ is
infinite.

On $\F_2^{\Z/m_\ell\Z}$ let $P\delta_t=\delta_{t+1}$. For $f$ depending
only on coordinates in $[-\ell,\ell]$ let
$D_\ell(f)\delta_t=f(T^ty_\ell)\delta_t$, where $f(T^ty_\ell)$ is the value
of $f$ at any point of $X$ that agrees with $T^ty_\ell$ on $[-\ell,\ell]$,
and put $\varphi_\ell(\sum_jf_ju^j)=\sum_jD_\ell(f_j)P^j$. Since
$PD_\ell(f)P^{-1}=D_\ell(f\circ T^{-1})$, for all $r,s\in R$ the identities
$\varphi_\ell(r+s)=\varphi_\ell(r)+\varphi_\ell(s)$ and
$\varphi_\ell(rs)=\varphi_\ell(r)\varphi_\ell(s)$ hold for large $\ell$. If
$r=\sum_jf_ju^j\ne0$, then for large $\ell$ some $D_\ell(f_j)\ne0$, and the
terms with $|j|<m_\ell/2$ have disjoint supports, so $\varphi_\ell(r)\ne0$.
For a finite $F\subseteq G$ and large $\ell$, applying $\varphi_\ell$ to the
entries of the elements of $F$ and of their inverses gives an injective
map $F\to\GL_{3m_\ell}(\F_2)$ that preserves the products staying in $F$.
So $G$ is LEF. LEF groups are sofic, and sofic groups are
hyperlinear~\cite[Example~4.5 and Theorem~3.3]{Pestov}.

\subsection*{Simplicity}
Let $1\ne N\trianglelefteq G$ and $1\ne g\in N$, and let $w\ge0$ bound
the absolute values of the exponents in the entries of $g$ and $g^{-1}$.
Call a clopen set $V$ small if $V\cap T^jV=\varnothing$ for
$0<|j|\le2w$ and every $f\circ T^i$, with $|i|\le w$ and $f$ a
coefficient of an entry of $g$ or $g^{-1}$, is constant on $V$.
Minimality and infiniteness imply that $T$ has no periodic points, so
every point has a small clopen neighborhood. By compactness every clopen
set is a finite disjoint union of small ones.

Some $h=e_{ij}(e_V)$ with $V$ small does not commute with $g$.
Otherwise $g$ commutes with $e_{ij}(e_V)$ for every clopen $V$,
by~\eqref{eq:elementary}. Taking $V=X$ gives $g=cI_3$ with
$c=\sum_jc_ju^j\in R$. Then $e_Vc-ce_V=\sum_jc_j(e_V-e_{T^jV})u^j$
vanishes for every clopen $V$. Choosing $V$ to contain $x$ but not
$T^{-j}x$, we get $c_j(x)=0$ for $j\ne0$. So $c\in\LC(X,\F_2)$, and
comparing constant coefficients in $cc^{-1}=1$ gives $c=1$ and $g=1$.

Fix such $h$ and $V$, put $d=3(2w+1)$, and for $|a|,|b|\le w$ put
$\epsilon_{ab}=u^ae_Vu^{-b}$. Since $e_Vu^ce_V=e_Ve_{T^cV}u^c$ vanishes
for $0<|c|\le2w$, we have
$\epsilon_{ab}\epsilon_{a'b'}=\delta_{ba'}\epsilon_{ab'}$. Index the
coordinates of $\F_2^d$ by pairs $(p,a)$ with $1\le p\le3$ and
$|a|\le w$. Then $\psi(E_{(p,a),(q,b)})=\epsilon_{ab}E_{pq}$ defines an
injective multiplicative linear map $\psi:M_d(\F_2)\to M_3(R)$, and
$A\mapsto I_3-\psi(I_d)+\psi(A)$ embeds $\GL_d(\F_2)$ in $\GL_3(R)$. Its
image $H$ lies in $G$, since transvections generate $\GL_d(\F_2)$. The
transvection between $(p,a)$ and $(q,b)$ maps to $e_{pq}(\epsilon_{ab})$
if $p\ne q$. If $p=q$, it is the commutator of the transvections between
$(p,a)$ and $(p',a)$ and between $(p',a)$ and $(p,b)$, where $p'\ne p$.

Put $k=[g,h]\in N\setminus\{1\}$. If $|a|,|b|\le w$ and $f,f'$ are
coefficients as above, then $fu^ae_Vf'u^b=u^a(f\circ T^a)e_Vf'u^b$ lies
in $\{0,\epsilon_{a,-b}\}$, since $f\circ T^a$ and $f'$ are constant on
$V$. The entries of $ghg^{-1}-I_3=g\,e_VE_{ij}\,g^{-1}$ are sums of such
products, and $h^{-1}=h=I_3+\epsilon_{00}E_{ij}$. So
$k-I_3=(ghg^{-1}-h)h$ and $k^{-1}-I_3=h(ghg^{-1}-h)$ lie in
$\psi(M_d(\F_2))$. As $\psi$ is injective and multiplicative, $k\in H$.
Then $N\cap H$ is a nontrivial normal subgroup of
$H\cong\GL_d(\F_2)=\PSL_d(\F_2)$, which is simple as $d\ge3$. So
$H\subseteq N$, and $e_{pq}(e_V)\in N$ for all $p\ne q$.

For distinct $p,q,l$ and $r,s\in R$, \eqref{eq:elementary} gives
$e_{pl}(re_V)=[e_{pq}(r),e_{ql}(e_V)]\in N$ and
$e_{pq}(re_Vs)=[e_{pl}(re_V),e_{lq}(s)]\in N$. By minimality and
compactness finitely many translates $T^aV$ cover $X$. Since
$e_{T^aV}=u^ae_Vu^{-a}$, the identity $1=1-\prod_a(1-e_{T^aV})$ writes
$1$, and so every element of $R$, as a sum of products $re_Vs$. So $N$
contains every $e_{pq}(r)$, and $N=G$.\hfill$\square$

Brown asks for an embedding into the unitary group of an
$\mathcal R^\omega$-embeddable McDuff factor, where $\mathcal R$ is the
hyperfinite $\mathrm{II}_1$ factor. As $G$ is hyperlinear, $L(G)$ embeds
in $\mathcal R^\omega$~\cite[Proposition~7.1]{Ozawa}. The infinite simple
group $G$ has infinite nontrivial conjugacy classes, so
$L(G)\mathbin{\bar\otimes}\mathcal R$ is a McDuff factor that embeds in
$\mathcal R^\omega$ and contains $G$ in its unitary group. With $3$
replaced by $n$ and $d=n(2w+1)$, the same proof shows that $\EL_n(R)$
has the properties of the theorem for every $n\ge3$.

\section{Word problems}

\begin{corollary}
The word problem of $G_X$ has the Turing degree of the language $L(X)$
of $X$. Every Turing degree occurs, so there are continuum many pairwise
nonisomorphic groups $G_X$.
\end{corollary}

\begin{proof}
Multiplying out a word in the generators gives a matrix with entries
$\sum_jf_ju^j$, each $f_j$ a table on words, using $uf=(f\circ T^{-1})u$.
The word is trivial if and only if the tables of its difference from
$I_3$ vanish on $L(X)$, so $L(X)$ computes the word problem. Conversely,
by~\eqref{eq:elementary} a word for $e_{12}$ of the cylinder indicator
$\prod_{t<n}u^{-t}e_{v_t}u^t$ is computable from $v=v_0\cdots v_{n-1}$,
and it is trivial if and only if $v\notin L(X)$. For derived topological
full groups, Grigorchuk and Medynets proved that the word problem is
decidable if and only if $L(X)$ is recursive~\cite[Theorem~1.1(3)]{GMpres}.

For irrational $\alpha\in(0,1)$ let $X_\alpha$ be the infinite minimal
Sturmian subshift, the closure of the codings $c(\theta)$,
$\theta\in[0,1)$, with $c(\theta)_t=1$ if and only if
$\theta+t\alpha\bmod1\ge1-\alpha$~\cite{MorseHedlund}. Its words of
length $n$ are the constant values of $c(\theta)_{[0,n)}$ on the arcs
between the points $-j\alpha\bmod1$, $0\le j\le n$, so $\alpha$ computes
$L(X_\alpha)$. The word $c(\theta)_{[0,n)}$ has
$\lfloor\theta+n\alpha\rfloor$ ones, within $1$ of $n\alpha$, so
$L(X_\alpha)$ computes $\alpha$. Every set $S\subseteq\mathbb N$ has the
degree of the irrational number $[0;1+\chi_S(0),1+\chi_S(1),\dots]$.
There are continuum many degrees, and the degree of the word problem is
an isomorphism invariant.
\end{proof}

\subsection*{Origin and authorship}
Claude (Anthropic) found the construction and original proofs under
the author's direction on September~12, 2026, and wrote the original
manuscript. Codex (OpenAI) shortened an earlier version, and Claude
revised this one.
The author is responsible for the final manuscript. Proof records are
in the \href{https://github.com/SauersML/group-approximation}{project repository}.

\begin{thebibliography}{99}

\bibitem{Brown}
N.~P. Brown, \emph{Tracial invariants, classification and
$\mathrm{II}_1$ factor representations of Popa algebras},
\href{https://arxiv.org/abs/math/0111286}{arXiv:math/0111286} (2001).

\bibitem{EJZ}
M.~Ershov and A.~Jaikin-Zapirain,
\emph{Property~\textup{(T)} for noncommutative universal lattices},
Invent. Math. \textbf{179} (2010), 303--347.
\doi{10.1007/s00222-009-0218-2}.

\bibitem{GM}
R.~Grigorchuk and K.~Medynets,
\emph{On algebraic properties of topological full groups},
Sb. Math. \textbf{205} (2014), 843--861.
\doi{10.1070/SM2014v205n06ABEH004400}.

\bibitem{GMpres}
R.~Grigorchuk and K.~Medynets,
\emph{Presentations of topological full groups by generators and
relations}, J. Algebra \textbf{500} (2018), 46--68.
\doi{10.1016/j.jalgebra.2016.10.027}.

\bibitem{Matui}
H.~Matui, \emph{Some remarks on topological full groups of Cantor
minimal systems}, Internat. J. Math. \textbf{17} (2006), 231--251.
\doi{10.1142/S0129167X06003448}.

\bibitem{MorseHedlund}
M.~Morse and G.~A. Hedlund,
\emph{Symbolic dynamics II. Sturmian trajectories},
Amer. J. Math. \textbf{62} (1940), 1--42.
\doi{10.2307/2371431}.

\bibitem{Ozawa}
N.~Ozawa, \emph{About the QWEP conjecture},
Internat. J. Math. \textbf{15} (2004), 501--530.
\doi{10.1142/S0129167X04002417}.

\bibitem{Pestov}
V.~G. Pestov, \emph{Hyperlinear and sofic groups: a brief guide},
Bull. Symbolic Logic \textbf{14} (2008), 449--480.
\doi{10.2178/bsl/1231081461}.

\bibitem{Thom}
A.~Thom, \emph{Examples of hyperlinear groups without factorization
property}, Groups Geom. Dyn. \textbf{4} (2010), 195--208.
\doi{10.4171/GGD/80}.

\bibitem{VershikGordon}
A.~M. Vershik and E.~I. Gordon,
\emph{Groups that are locally embeddable in the class of finite groups},
Algebra i Analiz \textbf{9} (1997), no.~1, 71--97; English transl.,
St. Petersburg Math. J. \textbf{9} (1998), no.~1, 49--67.

\end{thebibliography}
\end{document}
